% !TeX root=../main.tex
\chapter{آموزش سیاست انتخاب ویژگی و ارزیابی تجربی}
\label{chap:results}
\label{chap:policy-training}

\section{مقدمه}
در فصل سوم، چارچوب محاسبه توزیع برچسب همسایگی نمونه‌ها، مراکز کلاس‌ها و معیارهای فاصله و شباهت مورد نیاز برای طبقه‌بندی بررسی شد. هدف این فصل، صورت‌بندی روش یادگیری سیاست انتخاب ویژگی و ارزیابی تجربی آن است. در این صورت‌بندی، \emph{امتیاز طبقه‌بندی} که از شباهت بازنمایی نمونه به مراکز کلاس‌ها به دست می‌آید، از \emph{تابع هدف آموزش سیاست} که کیفیت تصمیم‌های حاصل از تغییر سیاست را می‌سنجد، متمایز است.

در آغاز، تولید مجموعه‌داده‌های مصنوعی با مدل مولد مشخص و معیار دقت بیزی معرفی می‌شود تا صحت محاسبات و سقف عملکرد مورد انتظار، مستقل از یادگیری سیاست، قابل بررسی باشد. سپس قاعده پیش‌بینی بر اساس بیشینه‌کردن شباهت نرمال‌شده جنسن--شنون با مراکز کلاس‌ها به کار گرفته می‌شود؛ خواص نظری این شباهت در فصل سوم بررسی شده‌اند. در مقابل، انتخاب تابع هدف، شیوه انتساب پاداش به تغییر احتمالات انتخاب ویژگی، روش به‌روزرسانی پارامترها و معیار توقف فرایند یادگیری، اجزای روش اجرایی این فصل هستند. ارزیابی نهایی نیز با استفاده از داده‌هایی انجام می‌شود که برچسب آن‌ها در آموزش سیاست و ساخت مراکز به کار نرفته است.


\section{تولید مجموعه‌داده‌های مصنوعی و تعریف حقیقت مبنا}
\label{sec:ch4-synthetic-data}
برای ارزیابی مستقل اجزای روش پیشنهادی، ابتدا مجموعه‌داده‌های مصنوعی با توزیع مولد معلوم تولید می‌کنیم. هدف از این انتخاب صرفاً دسترسی به نمونه‌های برچسب‌دار نیست، بلکه شناخت سازوکار تولید داده، ساختار آماری نمونه‌ها و رابطه احتمالاتی میان ویژگی‌ها و برچسب‌هاست. با کنترل عواملی مانند میزان هم‌پوشانی کلاس‌ها، قدرت تفکیک ویژگی‌ها، هم‌بستگی میان آن‌ها، فقدان مقادیر و عدم توازن کلاس‌ها، می‌توان رفتار روش پیشنهادی و نتایج آن را در شرایط گوناگون بررسی و اثر هر عامل را از سایر عوامل تفکیک کرد. مزیت دیگر این انتخاب آن است که علاوه بر برچسب هر نمونه، توزیع واقعی برچسب با توجه به ویژگی‌های مشاهده‌شده و توزیع برچسب مقصد در یک گام انتقال نیز قابل محاسبه است. بنابراین می‌توان خطای ناشی از تخمین گشت، کیفیت سیاست انتخاب ویژگی و عملکرد طبقه‌بند را از محدودیت ذاتی اطلاعات موجود در داده متمایز کرد. تولید داده و محاسبه کمیت‌های مرجع به صورت مستقل از الگوریتم آموزش انجام می‌شود.

\subsection{مدل مولد، هم‌بستگی ویژگی‌ها و فقدان داده}
\label{subsec:ch4-synthetic-model}
در نسخه پایه، $K$ کلاس و $d$ ویژگی عددی در نظر گرفته می‌شود. ابتدا برچسب یک نمونه از توزیع دسته‌ای با احتمال‌های پیشین $\pi_c>0$ انتخاب شده و سپس بردار کامل ویژگی‌ها از یک توزیع گاوسی چندمتغیره، مشروط بر کلاس، تولید می‌شود:
\begin{equation}\label{eq:ch4-synthetic-generative}
\boxed{
Y\sim\operatorname{Categorical}(\boldsymbol\pi),
\qquad
X\mid Y=c\sim\mathcal N(\boldsymbol\mu_c,\Sigma_c),
\qquad
\Sigma_c=D_c\,R\,D_c.
}
\end{equation}
در این رابطه $D_c=\operatorname{diag}(\sigma_{c,1},\ldots,\sigma_{c,d})$، و $R$ ماتریس هم‌بستگی متقارن، مثبت‌معین و با قطر واحد است که در پیاده‌سازی فعلی میان کلاس‌ها مشترک در نظر گرفته می‌شود. بدین ترتیب، میانگین و پراکندگی هر ویژگی می‌تواند در کلاس‌های مختلف متفاوت باشد و وابستگی میان ویژگی‌ها نیز به صورت صریح کنترل می‌شود. انتخاب $R=I$ حالت استقلال شرطی ویژگی‌ها را بازمی‌گرداند.

برخی ویژگی‌ها ممکن است توزیع حاشیه‌ای یکسان در تمام کلاس‌ها داشته باشند؛ چنین ویژگی‌هایی به‌تنهایی برای تشخیص برچسب اطلاع‌رسان نیستند و در اینجا «نویز حاشیه‌ای» نامیده می‌شوند. بااین‌حال، در صورت وجود هم‌بستگی با ویژگی‌های اطلاع‌رسان، همین ویژگی‌ها ممکن است در مشاهده مشترک نقش داشته باشند. بنابراین برای گزارش اهمیت تک‌ویژگی، از یک امتیاز توصیفی مبتنی بر میانگین وزن‌دار واگرایی KL متقارن میان توزیع‌های گاوسی حاشیه‌ای کلاس‌ها استفاده می‌کنیم؛ این امتیاز معادل اهمیت شرطی یا میزان تأثیر ویژگی بر دقت طبقه‌بندی نیست.

برای مدل‌کردن داده‌های ناقص، متغیر مشاهده‌پذیری $M_a\in\{0,1\}$ را برای هر ویژگی تعریف می‌کنیم. در مدل پایه، با نرخ فقدان $r_a$ داریم:
\begin{equation}\label{eq:ch4-synthetic-mask}
\mathrm P(M_a=0)=r_a,\qquad
M_a\ \perp\ (X,Y),
\qquad
\mathcal O=\{a:M_a=1\}.
\end{equation}
متغیرهای ماسک نیز در این نسخه مستقل از یکدیگر نمونه‌گیری می‌شوند؛ ازاین‌رو فقدان داده از نوع کاملاً تصادفی (\lr{MCAR}) است. مقدار ویژگی‌های خارج از مجموعه $\mathcal O$ از ورودی طبقه‌بند حذف و با ماسک مشخص می‌شود. نرخ‌های فقدان می‌توانند برای ویژگی‌های مختلف متفاوت باشند.

تمام پارامترهای مولد، شامل توزیع پیشین کلاس‌ها، بردارهای میانگین و انحراف معیار، ماتریس هم‌بستگی، نرخ فقدان هر ویژگی، اندازه مجموعه‌ها و بذر تصادفی، از فایل تنظیمات \lr{JSON} خوانده می‌شوند. در تنظیمات اولیه پیاده‌سازی، سه کلاس، پنج ویژگی و به‌ترتیب $300$، $90$ و $120$ نمونه برای آموزش، اعتبارسنجی و آزمون در نظر گرفته شده است. در این تنظیمات، ویژگی‌های با اندیس صفرمبنای $0$، $3$ و $1$ به‌ترتیب بیشترین تمایز حاشیه‌ای را دارند و ویژگی‌های $2$ و $4$ دارای توزیع حاشیه‌ای یکسان بین کلاس‌ها هستند. مقادیر یادشده تنها یک سناریوی قابل بازتولیدند و با تغییر فایل تنظیمات، بدون تغییر الگوریتم تولید، سناریوهای دیگری ساخته می‌شوند.

\subsection{محاسبه توزیع‌های واقعی و توزیع انتقال یک‌گامی}
\label{subsec:ch4-synthetic-oracles}
اگر $x_{\mathcal O}$ مقادیر مشاهده‌شده یک نمونه باشد، با معلوم‌بودن پارامترهای مولد، احتمال پسین واقعی کلاس به شکل زیر محاسبه می‌شود:
\begin{equation}\label{eq:ch4-synthetic-posterior}
\boxed{
\eta_c(x_{\mathcal O},M)
=\mathrm P(Y=c\mid X_{\mathcal O}=x_{\mathcal O},M)
=\frac{\pi_c\,\varphi_{|\mathcal O|}(x_{\mathcal O};
\boldsymbol\mu_{c,\mathcal O},\Sigma_{c,\mathcal O\mathcal O})}
{\sum_{r=1}^{K}\pi_r\,\varphi_{|\mathcal O|}(x_{\mathcal O};
\boldsymbol\mu_{r,\mathcal O},\Sigma_{r,\mathcal O\mathcal O})}.
}
\end{equation}
در اینجا $\varphi_{|\mathcal O|}$ چگالی گاوسی چندمتغیره روی ویژگی‌های مشاهده‌شده است. به دلیل مستقل‌بودن الگوی فقدان از برچسب و مقادیر ویژگی‌ها، احتمال ماسک در صورت و مخرج حذف می‌شود. اگر هیچ ویژگی‌ای مشاهده نشده باشد، این رابطه به $\eta_c=\pi_c$ کاهش می‌یابد. این توزیع، احتمال پسین \emph{واقعی} برچسب در مدل مصنوعی است و با توزیع برچسب همسایگی $\mathscr Y_i$ که از گشت تصادفی به دست می‌آید یکسان فرض نمی‌شود.

علاوه بر این، برای ارزیابی صحت انتقال عمودی، توزیع برچسب مقصد پس از \emph{یک گام} با انتخاب ویژگی $a$ را نیز محاسبه می‌کنیم. در مولد فعلی، وزن انتقال بر اساس کرنل
$\kappa_a(x_a,z_a)=\exp[-\gamma_a(x_a-z_a)^2]$، با $\gamma_a>0$، تعریف شده است. برای یک نمونه مقصد مستقل $(X',Y')$ از جامعه مولد داریم:
\begin{equation}\label{eq:ch4-synthetic-pop-q}
\begin{aligned}
Q_{x,a}^{\mathrm{pop}}(c)
&=\frac{\pi_c\,G_{c,a}(x_a)}
{\sum_{r=1}^{K}\pi_r\,G_{r,a}(x_a)},\\
G_{c,a}(x_a)
&=\frac{1}{\sqrt{1+2\gamma_a\sigma_{c,a}^{2}}}
\exp\left[
-\frac{\gamma_a(x_a-\mu_{c,a})^2}
{1+2\gamma_a\sigma_{c,a}^{2}}
\right].
\end{aligned}
\end{equation}
این رابطه از امید وزن کرنل نسبت به توزیع حاشیه‌ای ویژگی $a$ در هر کلاس به دست می‌آید؛ بنابراین هم‌بستگی بین ویژگی‌ها در $Q_{x,a}^{\mathrm{pop}}$ وارد نمی‌شود، هرچند در پسین مشترک رابطه~\eqref{eq:ch4-synthetic-posterior} مؤثر است. برای ویژگی‌ای که توزیع حاشیه‌ای آن در همه کلاس‌ها یکسان باشد، $Q_{x,a}^{\mathrm{pop}}(c)=\pi_c$. اگر ویژگی $a$ در نمونه آغازین موجود نباشد، این توزیع انتقال تعریف نمی‌شود.

همچنین $Q_{x,a}^{\mathrm{finite}}$ را از جمع وزن‌های کرنل روی نمونه‌های \emph{آموزش} که ویژگی $a$ را دارند، به تفکیک برچسب مقصد، به صورت دقیق محاسبه می‌کنیم. در محاسبه یک‌گامی نمونه‌های آموزش، خود نمونه از مجموعه مقصدها کنار گذاشته می‌شود؛ نمونه‌های اعتبارسنجی و آزمون فقط به داده‌های آموزش به‌عنوان مجموعه مرجع متصل می‌شوند. بدین ترتیب می‌توان توزیع یک‌گامی دقیق روی گراف محدود را از توزیع جامعه و نیز از تخمین حاصل از گشت تصادفی تفکیک کرد. محدودیت بازگشت به نمونه آغازین در گشت‌های چندگامی در پیاده‌سازی جداگانه گشت اعمال خواهد شد.

\subsection{دقت بیزی به عنوان سقف عملکرد مورد انتظار}
\label{subsec:ch4-bayes-ceiling}
ازآنجاکه توزیع پسین واقعی در رابطه~\eqref{eq:ch4-synthetic-posterior} معلوم است، می‌توان طبقه‌بند بیزی را برای اطلاعات مشاهده‌شده تعریف کرد:
\begin{equation}\label{eq:ch4-bayes-rule}
g_{\mathrm B}(x_{\mathcal O},M)
=\arg\max_{c=1,\ldots,K}\eta_c(x_{\mathcal O},M).
\end{equation}
این قاعده، تحت زیان صفر--یک، کمترین احتمال خطای مورد انتظار را در میان تمام طبقه‌بندهایی که به همین ویژگی‌ها و الگوی مشاهده‌پذیری دسترسی دارند به دست می‌آورد. بنابراین دقت بیزی جامعه برابر است با:
\begin{equation}\label{eq:ch4-bayes-accuracy}
\boxed{
\mathcal A_{\mathrm B}
=\mathbb E_{X_{\mathcal O},M}
\left[\max_{c}\eta_c(X_{\mathcal O},M)\right],
\qquad
\mathcal A(g)\le\mathcal A_{\mathrm B}.
}
\end{equation}
نامساوی بالا درباره \emph{امید ریاضی دقت} بر داده جدید با توزیع مولد یکسان و اطلاعات مشاهده‌شده یکسان است، نه درباره صحت تک‌تک پیش‌بینی‌ها. به دلیل هم‌پوشانی توزیع کلاس‌ها و حذف بخشی از ویژگی‌ها، ممکن است حتی طبقه‌بند بیزی نتواند تمام نمونه‌ها را درست تشخیص دهد. ازاین‌رو دقت بیزی \emph{تابع هدف آموزش سیاست نیست}؛ بلکه مرجع بالادستی عملکردی است که می‌توان از داده مشاهده‌شده انتظار داشت و برای تفسیر فاصله روش پیشنهادی تا بهترین تصمیم ممکن به کار می‌رود.

روی یک مجموعه آزمون محدود، دو برآورد متفاوت از عملکرد طبقه‌بند بیزی قابل محاسبه است. نخست، میانگین بیشترین احتمال پسین هر نمونه، که برآوردی از دقت بیزی مورد انتظار است؛ دوم، نسبت نمونه‌هایی که برچسب پیش‌بینی‌شده آن‌ها با برچسب واقعی یکسان است. این دو کمیت را به‌ترتیب چنین تعریف می‌کنیم:
\begin{equation}\label{eq:ch4-bayes-test-estimates}
\begin{aligned}
\widehat{\mathcal A}_{\mathrm B}^{\,\mathrm{prob}}
&=\frac{1}{n_{\mathrm{test}}}
\sum_{i=1}^{n_{\mathrm{test}}}
\max_c\eta_c(x_{i,\mathcal O_i},M_i),\\
\widehat{\mathcal A}_{\mathrm B}^{\,\mathrm{emp}}
&=\frac{1}{n_{\mathrm{test}}}
\sum_{i=1}^{n_{\mathrm{test}}}
\mathbf1_{\{g_{\mathrm B}(x_{i,\mathcal O_i},M_i)=y_i\}}.
\end{aligned}
\end{equation}
کمیت دوم تحت تأثیر تصادف در برچسب‌های مجموعه آزمون قرار دارد؛ بنابراین ممکن است روش دیگری روی یک مجموعه آزمون محدود تصادفاً دقت تجربی بالاتری داشته باشد، بدون آنکه کران بیزی در سطح جامعه نقض شده باشد.

همچنین باید میان دقت معمول و دقت متوازن تفاوت گذاشت. قاعده~\eqref{eq:ch4-bayes-rule} برای \emph{دقت معمول} بهینه است. اگر معیار نهایی دقت متوازن باشد، در مدل فعلی با پیشین‌های مثبت و فقدان مستقل از کلاس، قاعده بیزی متناظر به‌صورت زیر خواهد بود:
\begin{equation}\label{eq:ch4-bayes-balanced-rule}
g_{\mathrm B,bal}(x_{\mathcal O},M)
=\arg\max_c
\frac{\eta_c(x_{\mathcal O},M)}{\pi_c}.
\end{equation}
ازاین‌رو در گزارش نتایج، هر معیار عملکرد با مرجع بیزی متناظر خود مقایسه می‌شود.

\subsection{نقش داده‌های مصنوعی در مراحل بعدی}
\label{subsec:ch4-synthetic-usage}
برای هر سناریو، مجموعه‌های آموزش، اعتبارسنجی و آزمون به صورت مستقل و بازتولیدپذیر تولید می‌شوند. داده‌های قابل مشاهده برای الگوریتم شامل مقادیر موجود ویژگی‌ها، ماسک و برچسب‌های مجاز مجموعه آموزش است؛ توزیع پسین واقعی و توزیع‌های یک‌گامی مرجع در فایل‌هایی جداگانه ذخیره می‌شوند و تنها برای صحت‌سنجی و تحلیل نتایج مورد استفاده قرار می‌گیرند. به‌ویژه این توزیع‌های مرجع نباید ورودی روش آموزش سیاست یا پیش‌بینی نمونه آزمون باشند.

در گام نخست، خروجی انتقال یک‌گامی پیاده‌سازی با $Q^{\mathrm{finite}}$ و رفتار آن با $Q^{\mathrm{pop}}$ مقایسه می‌شود. پس از تأیید محاسبات گشت، سیاست یکنواخت مبنای ارزیابی قرار می‌گیرد و سپس تغییر سیاست، کیفیت بازنمایی و دقت طبقه‌بندی روی داده‌های اعتبارسنجی بررسی می‌شود. انتخاب نهایی تنظیمات و سیاست، بدون استفاده از برچسب‌های آزمون انجام خواهد شد؛ مجموعه آزمون فقط برای ارزیابی نهایی و مقایسه با سقف بیزی به کار می‌رود.

\section{پارامتردهی سیاست انتخاب ویژگی}
\label{sec:ch4-policy-param}
در ساختار گراف نمونه--ویژگی فصل سوم، هر گام شامل انتخاب یک ویژگی از نمونه جاری و سپس انتقال به نمونه‌ای دیگر بر اساس مقدار همان ویژگی است. سیاست انتخاب ویژگی فقط مرحله نخست را کنترل می‌کند؛ مرحله دوم با وزن ارتباط نمونه‌ها روی ویژگی انتخاب‌شده تعیین می‌شود. در این بخش، ابتدا خانواده عمومی سیاست‌های وابسته به وضعیت گشت را معرفی و سپس ساختار محدود و مستقل از برچسب را که مبنای آزمایش‌های این پژوهش است مشخص می‌کنیم.

\subsection{ساختار ماتریسی سیاست و ابعاد آن}
\label{subsec:ch4-policy-matrix}
فرض کنیم مجموعه ویژگی‌ها $A=\{1,\ldots,d\}$ باشد. انتخاب ویژگی در هر گام می‌تواند به ویژگی به‌کاررفته در انتقال قبلی وابسته باشد. علاوه بر $d$ وضعیت مربوط به ویژگی قبلی، برای شروع گشت که هنوز ویژگی پیشینی انتخاب نشده است، وضعیت متمایز $\bot$ را در نظر می‌گیریم. پارامترهای سیاست را در ماتریس امتیازهای خام زیر نگه می‌داریم:
\begin{equation}\label{eq:ch4-policy-logit-matrix}
\boxed{
Z=(z_{ab})_{a\in A\cup\{\bot\},\,b\in A}
\in\mathbb R^{(d+1)\times d}.
}
\end{equation}
درایه $z_{ab}$ امتیاز انتخاب ویژگی مقصد $b$ پس از ویژگی پیشین $a$ است و سطر $a=\bot$ امتیازهای نخستین انتخاب ویژگی را مشخص می‌کند. این درایه‌ها احتمال نیستند؛ ازاین‌رو لازم نیست نامنفی یا دارای مجموع سطری برابر یک باشند. در پیاده‌سازی، اندیس‌ها از صفر آغاز می‌شوند و سطر شماره $d$ به وضعیت شروع اختصاص دارد. برای مثال، با پنج ویژگی، ماتریس سیاست دارای شش سطر و پنج ستون است.

در این صورت‌بندی، یک ماتریس مشترک برای تمام نمونه‌ها و کلاس‌ها استفاده می‌شود. بنابراین امتیاز خام $z_{ab}$ به شناسه نمونه، مقدار ویژگی‌ها، برچسب نمونه آغازین یا برچسب نمونه جاری وابسته نیست. البته این استقلال فقط درباره \emph{پارامترهای خام} است؛ احتمال واقعی انتخاب ویژگی به محدودیت‌های نمونه جاری نیز وابسته خواهد بود.

\subsection{ویژگی‌های مجاز و محاسبه احتمال انتخاب}
\label{subsec:ch4-policy-masking}
برای گشتی که از نمونه $X_s$ آغاز شده و در نمونه $X_i$ قرار دارد، مجموعه ویژگی‌های قابل استفاده $B_s(X_i)$ مطابق رابطه~\eqref{eq:ch3-uniform-eligible} شامل ویژگی‌های موجود در نمونه جاری است که دست‌کم یک مقصد مجاز برای انتقال عمودی دارند. نمونه جاری و، در گام‌های بعدی، نمونه آغازین از مقصدهای مجاز حذف می‌شوند. احتمال انتخاب ویژگی $b$، پس از ویژگی قبلی $a$، با نرمال‌سازی نمایی تنها روی این مجموعه به دست می‌آید:
\begin{equation}\label{eq:ch4-policy-masked-softmax}
\boxed{
\pi_Z(b\mid i,a,s)=
\begin{cases}
\displaystyle
\frac{\exp(z_{ab})}
{\sum\limits_{u\in B_s(X_i)}\exp(z_{au})},
& b\in B_s(X_i),\\[7pt]
0,& b\notin B_s(X_i).
\end{cases}
}
\end{equation}
در گام اول $i=s$ و $a=\bot$ است. در هر گام بعدی، $a$ همان ویژگی انتخاب‌شده در گام پیشین خواهد بود. رابطه فوق در حالتی تعریف می‌شود که $B_s(X_i)$ ناتهی باشد؛ در پیاده‌سازی، نبود انتقال مجاز به‌عنوان خطای وضعیت گشت گزارش می‌شود و توزیع ساختگی جایگزین آن نمی‌گردد.

ازآنجاکه مجموعه مجاز می‌تواند بین نمونه‌ها متفاوت باشد، حتی با ماتریس ثابت $Z$، احتمال‌های نهایی انتخاب ویژگی برای دو نمونه الزاماً یکسان نیستند. افزون بر این، مقادیر واقعی ویژگی‌ها در احتمال \emph{انتقال به نمونه مقصد} اثر می‌گذارند، زیرا وزن انتقال عمودی $q_b^{(s)}(j\mid i)$ از فاصله مقادیر ویژگی $b$ محاسبه می‌شود. احتمال انتقال کامل برابر حاصل‌ضرب دو مرحله است:
\begin{equation}\label{eq:ch4-policy-complete-transition}
\mathrm P\bigl((X_j,b)\mid(X_i,a),X_s;Z\bigr)
=\pi_Z(b\mid i,a,s)\,q_b^{(s)}(j\mid i).
\end{equation}
بنابراین باید بین \emph{وابستگی سیاست انتخاب ویژگی به وضعیت گشت} و \emph{وابستگی وزن انتقال به مقدار ویژگی‌ها} تمایز گذاشت. همچنین پارامترهای $\beta$ و $\lambda$ که به‌ترتیب طول گشت و اهمیت مشاهدات در گشت را تعیین می‌کنند، جزو ماتریس $Z$ نیستند.

\subsection{امکان تعمیم به سیاست‌های وابسته به وضعیت و محدودیت مدل حاضر}
\label{subsec:ch4-policy-generalization}
در حالت عمومی می‌توان به‌جای نگهداری امتیاز ثابت برای هر زوج ویژگی، این امتیاز را خروجی تابع پارامتری در نظر گرفت:
\begin{equation}\label{eq:ch4-policy-general-score}
z_{ab}(s,i)=f_\theta
\bigl(a,b;\,x_s^{\mathrm{obs}},x_i^{\mathrm{obs}},
m_s,m_i,\mathcal I_{s,i}\bigr).
\end{equation}
در این رابطه $x^{\mathrm{obs}}$ مقادیر مشاهده‌شده ویژگی‌ها، $m$ ماسک وجود ویژگی‌ها و $\mathcal I_{s,i}$ سایر اطلاعات \emph{در دسترس} در وضعیت گشت است. چنین تابعی می‌تواند برای هر زوج ویژگی مبدأ و مقصد امتیاز متفاوتی، متناسب با مقدار ویژگی‌های نمونه جاری، ویژگی‌های نمونه آغازین یا اطلاعات دیگری از مسیر تولید کند. همچنین می‌توان سیاست را با چند ماتریس وابسته به کلاس یا یک مدل مشترک با ورودی کلاس شرطی کرد. در صورت استفاده از برچسب نمونه جاری، این برچسب فقط هنگامی قابل استفاده است که نمونه متعلق به مجموعه مرجع برچسب‌دار باشد؛ برچسب واقعی نمونه آزمون یا نمونه آغازین ناشناخته را نمی‌توان ورودی سیاست پیش‌بینی قرار داد. استفاده از فرضیه‌های مختلف برچسب آغازین نیز مستلزم تعریف یک روش استنتاج شرطی جداگانه است.

این خانواده عمومی از نظر توان نمایش انعطاف‌پذیرتر است، اما تعداد پارامترها، پیچیدگی محاسبه و نیاز به داده برای تخمین را افزایش می‌دهد. در آزمایش‌های پایه این پژوهش، از چنین تابع سازنده‌ای استفاده نمی‌کنیم و امتیازهای سیاست را به ماتریس مشترک و مستقل از مقدار و برچسب نمونه‌ها محدود می‌سازیم. این انتخاب، اثر احتمال‌های انتقال میان ویژگی‌ها را بدون ورود وابستگی‌های اضافی قابل بررسی می‌کند و امکان مقایسه مستقیم سیاست‌های ثابت و سیاست‌های قابل یادگیریِ هم‌ساختار را فراهم می‌آورد.

\subsection{سیاست‌های ثابت مبنا و وابستگی مراکز به سیاست}
\label{subsec:ch4-policy-baselines}
برای ارزیابی اولیه، دو ماتریس ثابت را در نظر می‌گیریم. در \emph{سیاست یکنواخت} تمام درایه‌های $Z$ صفرند؛ بنابراین هر ویژگی مجاز با احتمال $1/|B_s(X_i)|$ انتخاب می‌شود. در \emph{سیاست ثابت هدایت‌شده}، به ویژگی‌هایی که در سناریوی مصنوعی تمایز حاشیه‌ای بیشتری میان کلاس‌ها دارند امتیاز بزرگ‌تری داده می‌شود؛ همچنین امتیازها می‌توانند بسته به ویژگی پیشین متفاوت باشند. این سیاست بر اساس شناخت قبلی ساختار داده تنظیم شده است و نتیجه یادگیری خودکار محسوب نمی‌شود. برای مقایسه منصفانه این دو حالت، تنظیمات کرنل انتقال، طول گشت و وزن‌دهی مشاهده‌ها ثابت نگه داشته می‌شوند.

برای هر سیاست ثابت، ابتدا توزیع برچسب همسایگی نمونه‌های آموزش با همان سیاست محاسبه و سپس مراکز کلاس‌ها از میانگین این توزیع‌ها ساخته می‌شوند:
\begin{equation}\label{eq:ch4-policy-dependent-centers}
\mathcal V_c(Z)=
\frac{1}{|C_c|}
\sum_{i\in C_c}\mathscr Y_i(Z).
\end{equation}
در مرحله ارزیابی، ماتریس سیاست و مراکز ساخته‌شده ثابت می‌مانند. برای نمونه جدید تنها بازنمایی همسایگی با همان سیاست محاسبه می‌شود و برچسب از مقایسه آن با مراکز ذخیره‌شده به دست می‌آید؛ نه سیاست و نه مراکز بر اساس نمونه آزمون به‌روزرسانی نمی‌شوند. به این ترتیب، نقش سیاست در \emph{ساخت مراکز} از نقش آن در \emph{جست‌وجو و پیش‌بینی} جدا ولی سازگار باقی می‌ماند.


\section{ساخت بازنمایی و قاعده پیش‌بینی}
\label{sec:ch4-inference}
\subsection{تخمین توزیع همسایگی با گشت‌های تصادفی}
\subsection{محاسبه مراکز کلاس‌ها و امتیاز شباهت}
\subsection{پیش‌بینی برچسب نمونه‌های جدید}

\section{تابع امتیاز آموزش و بهینه‌سازی سیاست}
\label{sec:ch4-policy-objective}
\subsection{حاشیه تصمیم و خطای طبقه‌بندی متوازن}
\subsection{تعریف پاداش و ارتباط آن با تغییر سیاست}
\subsection{روش به‌روزرسانی پارامترها و معیار توقف}

\section{طرح آزمایش و روش ارزیابی}
\label{sec:ch4-evaluation-protocol}
\subsection{مجموعه‌داده‌ها، پیش‌پردازش و نحوه تفکیک آموزش و آزمون}
\subsection{سیاست یکنواخت و سایر روش‌های مبنا}
\subsection{معیارهای ارزیابی و تکرار آزمایش‌ها}

\section{نتایج تجربی و تحلیل}
\label{sec:ch4-results}
\subsection{دقت طبقه‌بندی پیش و پس از آموزش سیاست}
\subsection{تحلیل حاشیه تصمیم، جدایی مراکز و پراکندگی درون‌کلاسی}
\subsection{حساسیت به پارامترها و هزینه محاسباتی}
